\documentclass[11pt,a4paper]{article}

\usepackage[T1]{fontenc}
\usepackage[utf8]{inputenc}
\usepackage{lmodern}
\usepackage[english]{babel}
\usepackage[a4paper,margin=2.2cm]{geometry}
\usepackage{amsmath,amssymb}
\usepackage{xcolor}
\usepackage{enumitem}

\definecolor{accent}{HTML}{0F4C5C}

\setlength{\parskip}{0.5em}
\setlength{\parindent}{0pt}

\newcommand{\slide}[1]{\section*{#1}}
\newcommand{\points}[1]{{\color{accent}\textbf{Slide content:}}\\#1}
\newcommand{\speech}[1]{{\color{accent}\textbf{What to say:}}\\#1}

\title{Superconducting Devices\\
\large From Strong Magnetic Fields to Ultra-Sensitive Sensing}
\author{Filippo Vicidomini\\University Ca' Foscari of Venice}
\date{April 2026}

\begin{document}
\maketitle

\section*{Overview}

\textbf{Core idea:}
\begin{quote}
The same superconducting physics leads to two technological routes:
strong-field magnets and ultra-sensitive SQUID sensors.
\end{quote}

\begin{itemize}
\item minimal theory
\item SQUID principle
\item three applications
\end{itemize}

\newpage

%------------------------------------------------

\slide{1. Title}

\points{
Title, name, course. Clean background.
}

\speech{
Superconductivity is often introduced as a beautiful but quite abstract topic in condensed matter physics. In this talk, I want to show the opposite perspective: the same physical principles can be turned into very concrete technologies.

In particular, superconductivity opens two complementary routes. On one side, it allows us to generate very strong magnetic fields with very low dissipation. On the other side, it allows us to detect extremely weak magnetic signals with extraordinary precision.

So the main message of the talk is simple: one underlying physics, but two very different device families, leading to very different applications.
}

%------------------------------------------------

\slide{2. One Physics, Two Device Families}

\points{
Superconductivity → High-field magnets / SQUID sensors

Applications:
\begin{itemize}
\item Water treatment
\item MEG
\item Buried-metal detection
\end{itemize}
}

\speech{
The structure of the talk is intentionally simple. I will begin with the essential physical properties of superconductors, then move to the Josephson effect, which is the key ingredient for superconducting electronics.

From there, I will show how the same physics leads to two device families. The first is based on producing strong magnetic fields, mainly through superconducting magnets. The second is based on ultra-sensitive magnetic sensing, mainly through SQUIDs.

Finally, I will connect these ideas to three concrete applications: water treatment, magnetoencephalography, and buried-metal detection.
}

%------------------------------------------------

\slide{3. What is a Superconductor?}

\points{
\begin{itemize}
\item Zero electrical resistance
\item Meissner effect
\item Type-II vortices
\end{itemize}
}

\speech{
A superconductor is not just a material with very low resistance. It is a distinct quantum state of matter that appears below a critical temperature.

The first famous property is zero electrical resistance, meaning that an electric current can flow without dissipation. But this alone is not the full story.

The defining feature is the Meissner effect: magnetic fields are expelled from the interior of the material. This tells us that superconductivity is not simply perfect conductivity, but a new thermodynamic phase with its own electromagnetic behavior.

At the microscopic level, electrons form bound pairs, called Cooper pairs, and these pairs condense into a single collective quantum state. So instead of thinking about many independent electrons, we describe the system with one macroscopic wavefunction.

In many practical materials, especially type-II superconductors, magnetic flux can still enter in the form of quantized vortices when the external field is high enough. This is extremely important technologically, because it allows superconductivity to survive even in strong magnetic fields.

So the key point to remember is this: superconductivity combines zero resistance, magnetic-field expulsion, and macroscopic quantum coherence in the same physical state.
}

%------------------------------------------------

\slide{4. Macroscopic Phase and Josephson Effect}

\points{
\[
\psi = |\psi| e^{i\theta}
\]

\[
I_s = I_c \sin \varphi
\]

(Optional)
\[
\frac{d\varphi}{dt} = \frac{2eV}{\hbar}
\]
}

\speech{
The superconducting state can be described by a macroscopic wavefunction, written as $\psi = |\psi| e^{i\theta}$. The important quantity here is the phase $\theta$.

In a normal conductor, phase is not something we can use at the macroscopic level. In a superconductor, instead, the phase remains coherent over large distances and becomes physically meaningful.

This leads directly to the Josephson effect. If two superconductors are separated by a very thin insulating barrier, Cooper pairs can tunnel from one side to the other.

The first Josephson relation tells us that the supercurrent depends on the phase difference: $I_s = I_c \sin \varphi$. So a quantum variable, the phase difference, directly controls a measurable current.

The second relation says that if a voltage is applied, the phase difference changes in time, $\frac{d\varphi}{dt} = \frac{2eV}{\hbar}$. So current is linked to phase, and voltage is also linked to phase.

This is the conceptual bridge we need for the rest of the talk. Superconductivity is not only about lossless current transport, but also about controlling a macroscopic quantum phase. And once phase becomes accessible, we can build devices based on interference.
}

%------------------------------------------------

\slide{5. From Josephson Junction to SQUID}

\points{
\begin{itemize}
\item Superconducting loop
\item Two Josephson junctions
\item Magnetic flux modifies interference
\end{itemize}

\[
\Delta \Phi \rightarrow \Delta V
\]
}

\speech{
If we place two Josephson junctions in a superconducting loop, we obtain a SQUID: a Superconducting Quantum Interference Device.

The key idea is that the superconducting phase around the loop cannot vary arbitrarily. Because of flux quantization, the total phase change around the loop is constrained by the magnetic flux threading the loop.

So when an external magnetic flux is applied, it changes the phase configuration of the device. The two junctions therefore interfere with each other, and the total critical current becomes a periodic function of magnetic flux.

In practice, the device is biased in a regime where this dependence is very steep. Then even a tiny change in magnetic flux produces a measurable change in the electrical output, usually a voltage.

So the chain is the following: magnetic flux changes the phase, the phase changes the interference, and the interference changes the electrical signal.

This is why the SQUID is such a powerful sensor: it converts extremely small magnetic signals into measurable voltages through quantum interference.
}

%------------------------------------------------

\slide{6. Why SQUIDs Matter}

\points{
\begin{itemize}
\item Ultra-high sensitivity
\item Detects fields far below Earth's background
\item Passive sensing
\item Works where other sensors fail
\end{itemize}
}

\speech{
The main reason SQUIDs are so important is their extraordinary sensitivity. They are among the most sensitive magnetometers ever developed.

Under suitable conditions, they can detect magnetic fields down to the femtotesla range. To appreciate how small this is, the Earth's magnetic field is of the order of tens of microtesla, so SQUIDs can detect signals many orders of magnitude weaker than the magnetic background around us.

This makes them essential in applications where the signal is extremely weak, where the measurement must be non-invasive, or where low-frequency performance is crucial. In many of these situations, conventional magnetic sensors are simply not sensitive enough.
}

%------------------------------------------------

%------------------------------------------------

\slide{7. Application: MEG}

\points{
\begin{itemize}
\item Neuronal currents → magnetic fields
\item Signals $\sim 10$--$100$ fT
\item Non-invasive measurement
\item SQUID array ($\sim$300 sensors)
\end{itemize}
}

\speech{
Magnetoencephalography, or MEG, is one of the clearest real-world applications of SQUID technology.

The signal originates from neuronal currents in the brain, mainly due to synchronized activity of pyramidal neurons. These currents generate magnetic fields that extend outside the head.

However, these signals are extremely weak, typically in the range of 10 to 100 femtotesla. To give a sense of scale, they are about a billion times smaller than the Earth's magnetic field.

This is exactly where SQUIDs become essential. Their extreme sensitivity allows us to detect these signals in a completely non-invasive way.

In practice, MEG systems use large arrays of SQUID sensors, often around 300 channels, placed inside a cryogenic helmet surrounding the head. The system records brain activity with millisecond temporal resolution, making it ideal for studying fast neural dynamics.
}

%------------------------------------------------

\slide{7b. MEG System and Advantages}

\points{
\begin{itemize}
\item Cryogenic helmet (dewar)
\item ms temporal resolution
\item Less distortion than EEG
\item Better source localization
\end{itemize}
}

\speech{
From a system perspective, MEG relies on a cryogenic environment. The SQUID sensors must be kept at very low temperature, typically using liquid helium, inside what is called a dewar.

One of the main strengths of MEG is its excellent temporal resolution, on the order of milliseconds, which allows us to follow brain activity in real time.

Compared to EEG, magnetic fields are less affected by the skull and scalp. This means that the measured signals are less distorted, and source localization is more direct.

So overall, MEG combines non-invasive measurement, high temporal resolution, and relatively accurate spatial localization, making it a powerful tool both in neuroscience research and in clinical applications such as epilepsy diagnosis.
}

%------------------------------------------------

\slide{8. Application: Buried-Metal Detection}

\points{
\begin{itemize}
\item Excitation field
\item Induced currents in target
\item Weak secondary field
\item SQUID detection
\end{itemize}
}

\speech{
This application belongs to the sensing branch, but the mechanism is different from MEG. Here the key idea is electromagnetic induction.

A primary coil generates an alternating magnetic field that penetrates the ground. When this field encounters a buried metallic object, it induces circulating currents inside the target, called eddy currents.

These eddy currents generate a secondary magnetic field. So the buried object does not simply block the field: it produces its own magnetic response, and this response carries information about the presence, size, and shape of the object.

At this point, the natural question is: why are SQUIDs useful here, instead of an ordinary metal detector?

The answer is that the secondary field is extremely weak compared both to the excitation field and to the environmental magnetic background. So the real challenge is not generating the signal, but detecting it reliably.
}

%------------------------------------------------

\slide{8b. Buried-Metal Detection: Why SQUIDs?}

\points{
\begin{itemize}
\item Extremely weak secondary field
\item High sensitivity in the fT range
\item Low-frequency operation \(\rightarrow\) deeper penetration
\item UXO and mine detection
\end{itemize}
}

\speech{
This is exactly where SQUID sensors become valuable.

Because of their extremely high sensitivity, they can detect the weak secondary magnetic field produced by the buried target, even in a noisy environment.

Another important advantage is that SQUIDs can operate at low frequencies. This is useful because lower frequencies penetrate deeper into the ground, so the system can detect objects buried at greater depths.

This is why SQUID-based systems are interesting not only for generic metal detection, but also for more demanding applications such as unexploded ordnance detection.

In that context, the goal is not just to say that some metal is present. The system should help distinguish dangerous buried objects from harmless metallic debris, which is one of the reasons why high sensitivity and better signal characterization matter.
}

%------------------------------------------------

\slide{9. Application: Water Treatment}

\points{
\begin{itemize}
\item Magnetic tagging of pollutants
\item HGMS as a magnetic filter
\item Strong local gradients: need high $B \nabla B$
\item Removal efficiency $>$95\%
\end{itemize}
}

\speech{
This application shows the strong-field branch of superconductivity.

The starting point is that most pollutants are not magnetic, so we cannot remove them directly using a magnetic field.

To solve this, we first attach pollutants to magnetic particles, typically iron oxide particles or nanoparticles. In this way, contaminants become magnetically responsive.

The water is then sent through a high-gradient magnetic separation system, or HGMS. In practical terms, an HGMS system includes an excitation coil, a magnetic yoke, and a separation region containing a ferromagnetic matrix. In some industrial designs, there is also a rotating ring or pulsating mechanism that helps continuous operation and reduces clogging.

Under the action of the water flow, the suspension enters the separator. Inside the magnetic field, the tagged particles are attracted toward the magnetized matrix and form magnetic aggregates, while the non-magnetic part of the flow continues forward.

As the separator rotates or as the flow is periodically flushed, the captured material is carried to a lower-field region and removed. So the process is not only capture, but also continuous collection and discharge of the concentrated waste.

Because of the ferromagnetic structure, very large local field gradients are created. These gradients generate the magnetic force that pulls the tagged particles toward the matrix, while cleaner water continues to flow.

In many practical implementations, removal efficiencies above 95\% are reported, even for very small particles, which makes this technique highly effective for water purification.
}

%------------------------------------------------

\slide{9b. Water Treatment: Role of Superconductivity}

\points{
\begin{itemize}
\item Need strong magnetic field
\item Create high gradients
\item Continuous separation
\item Energy-efficient at scale
\end{itemize}
}

\speech{
The key technological role of superconductivity here is to generate strong and stable magnetic fields.

These strong fields magnetize the iron yoke and the ferromagnetic matrix inside the separator, creating the high-gradient conditions needed for capture. So the important quantity is not only the field itself, but also its spatial variation.

This is why HGMS is more than simply putting water near a magnet. The device must combine field generation, flow control, and an engineered magnetic structure that concentrates the gradients where separation takes place.

Superconducting magnets are especially useful because they can provide these strong fields continuously with low energy dissipation, which is very important for large-scale systems.

So, differently from MEG, where superconductivity is used for sensitivity, here it is used to create the magnetic environment needed to physically separate contaminants from water.
}

%------------------------------------------------

\slide{10. Comparison}

\points{
\begin{tabular}{lccc}
Application & Main need & Device \\
\hline
MEG & Extreme sensitivity & SQUID \\
Metal detection & Weak-signal detection & SQUID \\
Water treatment & Strong magnetic force & Magnet \\
\end{tabular}
}

\speech{
Although these three applications look very different, they fit into one simple framework.

In MEG and buried-metal detection, the central requirement is sensitivity to extremely weak magnetic signals. In both cases, the enabling device is the SQUID.

In water treatment, instead, the important quantity is magnetic force, which depends on strong fields and strong gradients. There the enabling technology is the superconducting magnet.

So the underlying physics is shared, but the engineering objective changes. In one case we use superconductivity to measure tiny magnetic signals. In the other, we use it to generate the magnetic conditions needed for separation.
}

%------------------------------------------------

\slide{11. Limits and Outlook}

\points{
\textbf{Limits}
\begin{itemize}
\item Cryogenics
\item Noise and shielding
\item System complexity
\end{itemize}

\textbf{Outlook}
\begin{itemize}
\item High-$T_c$ devices
\item Simpler cooling
\item Better integration
\end{itemize}
}

\speech{
Despite their advantages, superconducting devices also come with important limitations.

The first is cryogenics. In most implementations, we still need low temperatures, and this increases cost, complexity, and maintenance requirements.

In sensing applications, another major issue is environmental noise. Since the signals can be extremely weak, shielding and careful system design become essential.

There is also a broader engineering challenge: the superconducting element alone is never the whole system. We also need readout electronics, mechanical integration, thermal management, and reliability.

Looking ahead, much of the progress in the field is driven by high-temperature superconductors, simpler cooling solutions, and more compact and robust devices. So the long-term trend is to keep the unique performance of superconducting systems while reducing the practical barriers to their use.
}

%------------------------------------------------

\slide{12. Conclusion}

\points{
Superconductivity →
\begin{itemize}
\item Create magnetic fields
\item Detect magnetic fields
\end{itemize}

Applications:
\begin{itemize}
\item Water treatment
\item MEG
\item Metal detection
\end{itemize}
}

\speech{
In conclusion, superconductivity is not only an elegant quantum phenomenon. It is also a technological platform.

The same physical principles allow us to do two things extremely well: generate strong magnetic fields and detect extremely weak ones.

The examples discussed in this talk show these two routes clearly. Water treatment belongs to the strong-field branch, while MEG and buried-metal detection belong to the ultra-sensitive sensing branch.

So the final message is that superconductivity gives us a way to both shape and measure the magnetic world with exceptional precision.
}

\end{document} 